\section{Introduction}

The minimum-variance portfolio is the solution to a quadratic programme whose
inputs are a covariance estimate and a set of weight constraints. In the large-$n$
regime of institutional equity portfolios ($n/T \approx 0.4$ for a five-year daily
window), the sample covariance $\hat\Sigma = X^\top X/T$ is nearly singular,
with condition number $\kappa(\hat\Sigma) \approx 106{,}000$ on S\&P~500 data.
This ill-conditioning makes the portfolio optimisation problem numerically fragile.

Random matrix theory (RMT) attributes most of this ill-conditioning to noise in
the \emph{correlation} structure: applied to the sample correlation
$C = D^{-1/2}\hat\Sigma\,D^{-1/2}$ ($D = \mathrm{diag}(\hat\Sigma)$), the
Marchenko--Pastur (MP) law identifies correlation eigenvalues below the bulk edge
$\hat\lambda_+ = (1+\sqrt{n/T})^2$ as noise and those above as genuine factor
structure. The ``constant residual eigenvalue'' (CRE) cleaning
method~\cite{laloux1999,bun2017,bouchaud2003} replaces the noise eigenvalues with
their average $\bar\mu$ (preserving $\mathrm{tr}(C) = n$) and produces the cleaned
correlation
\[
  C_0 = \bar\mu\, I + U_k(\Lambda_k - \bar\mu I)U_k^\top,
\]
where $(U_k, \Lambda_k)$ are the $k$ correlation eigenpairs above $\hat\lambda_+$.
On S\&P~500 data, $k = 17$ and $\kappa(C_0) \approx 313$, roughly $290$ times
better conditioned than the sample correlation ($\kappa(C) \approx 91{,}000$).

Implementations of CRE cleaning in the public domain typically form the cleaned
\emph{covariance} as a dense $n \times n$ matrix and pass it to a general-purpose
quadratic-program solver. This paper exploits the
scalar-identity-plus-low-rank structure of $C_0$, which renders the matrix inverse
available in closed form via the Woodbury identity, without assembling either the
$n \times n$ covariance or the $n_a \times n_a$ active-set restriction. Solving in
the standardised coordinates $y = D^{1/2}w$ (exact, since
$w \geq 0 \Leftrightarrow y \geq 0$) poses the minimum-variance problem against
$C_0$ directly, so the solver inherits the cleaned correlation's conditioning. The
result is a direct active-set portfolio solver substantially faster than the
KKT-Cholesky approach of assembling the restricted system and factoring.
Exploiting factor structure in portfolio quadratic programmes via the
Sherman--Morrison--Woodbury identity goes back at least to
Perold~\cite{perold1984}; the contribution here is the combination with
RMT-determined structure, randomised preprocessing, and an end-to-end
complexity and accuracy account.

\paragraph{Contributions}
This paper makes three algorithmic contributions and one empirical
contribution.

\emph{First, a direct Woodbury portfolio solver.} The scalar-identity-plus-low-rank
form of $C_0$ renders the active-set inner solve available in
closed form via the Woodbury identity, at $O(n_a k^2 + k^3)$ per outer step in
place of the $O(n_a^2 k + n_a^3)$ assemble-and-factor cost, without forming
either the $n \times n$ covariance or its $n_a \times n_a$ active-set
restriction (Section~\ref{sec:algorithm}). The $p$ linear balance constraints
$Bw=c$ (the budget $\mathbf{1}^\top w=1$ as the leading case) are eliminated
analytically through a $p \times p$ Schur complement. We derive the exact crossover
condition against KKT-Cholesky (Proposition~\ref{prop:crossover}): the Woodbury
inner solve is cheaper per outer step precisely when $k < n_a$, a condition
satisfied throughout the long-only equity regime ($k/n_a$ between roughly $0.2$
and $0.7$ in our experiments), and we account for the total pipeline cost over
a sweep of $S$ related solves (Proposition~\ref{prop:pipeline}), under which the
$O(nTk)$ preprocessing amortises to $O(nTk/S)$ per solve.

\emph{Second, randomised-SVD preprocessing.} The naive path forms $\hat\Sigma$
and takes a dense eigendecomposition, at $O(n^2 T + n^3)$ cost and $O(n^2)$
storage. A randomised truncated SVD applied directly to $X$ computes the
top-$k$ eigenpairs of $\hat\Sigma$ in $O(nTk)$ time and $O(nk)$ storage
(Section~\ref{sec:preprocessing}); the measured portfolio impact of the
approximation is below three basis points in any weight. For daily rolling
windows we additionally implement an incremental rank-two update of the tracked
eigenpairs that avoids re-sketching altogether.

\emph{Third, the scope of the structural trick}
(Remarks~\ref{rem:rie} and~\ref{rem:ridge}): ridge terms, diagonal loadings,
externally specified factor models, and linear equality constraints preserve
the Woodbury structure, whereas statistically richer cleanings (nonlinear
shrinkage and general rotationally invariant estimators) do not.

\emph{Fourth, empirical validation.} We compare Woodbury directly against
KKT-Cholesky on the same RMT problem (Section~\ref{sec:results}), with a KKT
optimality certificate for the returned weights and an exact critical-line sweep
built on the same Woodbury kernel that traces every frontier breakpoint
(Section~\ref{ssec:cla}). A $k$-sensitivity
audit shows the MP threshold is not a sensitive hyperparameter
(Section~\ref{sec:ksensitivity}), and rolling out-of-sample backtests on two
universes (S\&P~500 and FTSE~100) show the cleaned estimator carries no
measurable statistical penalty relative to standard shrinkage
(Section~\ref{ssec:oos}).

\paragraph{Relation to existing work}
Laloux et al.~\cite{laloux1999} and Plerou et al.~\cite{plerou1999} brought the
Marchenko--Pastur threshold to financial correlations; Bouchaud and
Potters~\cite{bouchaud2003} and Bun, Bouchaud and Potters~\cite{bun2017} develop
the theory, including the statistically superior rotationally invariant estimators
of Ledoit and P\'ech\'e~\cite{ledoitpeche2011} and the nonlinear shrinkage of
Ledoit and Wolf~\cite{ledoit2012,ledoit2020}. Those modify the full spectrum and
are not Woodbury-compatible; we work with CRE because clipping is the member of the
family with exploitable structure (Remark~\ref{rem:rie}), popularised for
practitioners by L\'opez de Prado~\cite{lopezdeprado2018}. On the optimisation
side, Perold~\cite{perold1984} exploited diagonal-plus-low-rank factor structure in
a portfolio QP, the critical line algorithm of
Markowitz~\cite{markowitz1956,niedermayer2010,cvxcla} traces the exact long-only
frontier by active-set pivoting, and the randomised SVD is due to Halko, Martinsson
and Tropp~\cite{halko2011}. Our contribution is the pairing of RMT-determined
low-rank structure with the active-set solve, the randomised preprocessing with an
explicit pipeline complexity, and the measured end-to-end comparison.

This paper is self-contained: Algorithm~\ref{alg:woodbury} gives the complete
primal-dual active-set solver including the Woodbury inner solve, and
Section~\ref{sec:preprocessing} develops the rSVD preprocessing independently.

\paragraph{Outline}
Section~\ref{sec:rmt} constructs the cleaned estimator.
Section~\ref{sec:problem} states the long-only problem.
Section~\ref{sec:algorithm} develops the active-set reduction, the Woodbury
portfolio algorithm, and the crossover analysis. Section~\ref{sec:preprocessing} covers
randomised SVD preprocessing. Section~\ref{sec:results} reports timing, scaling,
threshold sensitivity (Section~\ref{sec:ksensitivity}), and out-of-sample
results. Section~\ref{sec:conclusions} concludes.
